\section{서론}

\subsection{배경: 방산 자율체계의 사이버 위협 지형}

무인항공기 (UAV), 무인지상차량 (UGV), 그리고 이들의 원격 지휘통제 (C2) 메시지를 전달하는 위성통신 (SATCOM) 기반 비가시선 (BLOS) C2 링크는 현대 방산 자율체계의 핵심 기반이다. 
%
방산 무인체계의 운용 환경은 민간용 드론과 본질적으로 다르다. 
%
가장 중요한 차이는 적의 감시, 교란, 기만, 침투 가능성을 기본 조건으로 삼는 적대적 작전환경 (Contested operational environment)에 있다. 본 연구는 이러한 적대성을 세 가지 방산 특수성으로 구조화한다.

\textbf{첫째}, BLOS 운용은 경로상 노출을 구조적으로 만든다. 
%
방산 정보·감시·정찰 (Intelligence, Surveillance, and Reconnaissance, ISR) 및 타격 UAV와 원거리 UGV는 가시선 밖에서 임무를 수행하므로, C2 링크가 위성 중계 또는 다중 홉 무선주파수 (Radio Frequency, RF) 경로에 의존한다. 이때 임차 트랜스폰더, bent-pipe 중계기, 지상 게이트웨이는 모두 관측, 지연, 손실, 변조가 가능한 개입 지점이 된다.
%
\textbf{둘째}, 글로벌 위성항법시스템 거부는 예외 상황이 아니라 상시 위협 조건이다. 현대 전장에서는 GPS 재밍과 기만이 항법 체계를 직접 겨냥한다. 
%
무인체계가 GNSS 단일 소스에 크게 의존할수록 위치·항법·시각 (positioning, navigation, and timing, PNT) 정보의 무결성은 단순한 센서 신뢰성 문제가 아니라 임무 지속성과 직결되는 보안 문제가 된다.
%
\textbf{셋째}, 레거시 MAVLink 운용 구성은 명령 출처성과 메시지 무결성을 충분히 보장하지 못한다. 
%
마이크로 항공기 링크 (MAVLink) v1은 인증, 서명, 암호화를 기본적으로 제공하지 않는다. MAVLink v2는 메시지 서명을 지원하지만, 서명을 명시적으로 활성화하지 않는 운용 구성에서는 RC 오버라이드, 명령 주입, 미션 변조, 파라미터 변조에 대한 프로토콜 수준의 출처 검증을 기대하기 어렵다. 또한 MAVLink v2 서명은 메시지 출처와 무결성을 검증하는 기능이지, 메시지 내용 자체를 암호화하는 기능은 아니다.
%
이를 요약하면, 본 연구의 위협 모델은 인증이 부재한 제어 프로토콜의 약점과 BLOS 및 전자전 (electronic warfare, EW) 환경에서 발생하는 경로상 접근 가능성의 결합으로 정의된다. 즉, 핵심 공격면은 단순한 민수용 드론 해킹이 아니라, 위성/BLOS C2 운용 구조에서 발생하는 메시지 출처성, 무결성, 재전송 저항성의 결여다.

공개 문헌과 실제 사건 사례도 이러한 위협 지형을 뒷받침한다. 이란의 RQ-170 나포 사건은 C2 거부와 GNSS 기만 가능성이 함께 논의되는 대표 사례로 인용된다\cite{humphreys_testimony_2012}. 우크라이나 전장의 전자전 사례는 재밍, 기만, C2 방해가 무인체계 운용에 지속적으로 영향을 주는 환경을 보여준다\cite{rusi_russian_airwar_2022}. CVE-2020-10283은 MAVLink 버전 협상 및 인증 우회 가능성을 기록하며, sysid 사칭과 명령 주입 공격군의 프로토콜적 배경을 제공한다\cite{cve202010283}. 한편 Psiaki와 Humphreys의 무결 (Seamless) GPS 기만 탈취 연구\cite{psiaki_humphreys_2016}, Son 등의 음향 공진 기반 MEMS 자이로 공격\cite{son2015rocking}, Trippel 등의 WALNUT 음향 MEMS 가속도계 기만 연구\cite{trippel2017walnut}는 현실적 참조 위협 (Reference threat)으로 인용하되, 직접 구현 대상에서는 제외한다. 이 구분은 본 연구의 정직성 원칙과 직접 연결된다.

\subsection{문제 정의: 출처 모호성 (provenance ambiguity)}
% 방산 자율체계 제어 보안의 근본 난제는 \textbf{출처 모호성 (Provenance ambiguity)}으로 압축된다. 자율 비행·주행 중 오토파일럿은 정당한 명령과 경로상 주입된 명령을 작동기 수준 (Actuator level)에서 구별할 수 없다. 위조된 단일 채널 RC 오버라이드 프레임은 작동기 관점에서 정당한 조종 명령과 바이트 단위로 동일하다. 따라서 사후 탐지 (Detection)만으로는 이 공격 클래스를 원천 차단할 수 없으며, 여기서 두 방어 설계 원리가 직접 도출된다.

% \begin{enumerate}[leftmargin=*]
%     \item \textbf{1차 방어는 인증이어야 한다.} 출처 모호성을 근본적으로 해소하는 유일한 방법은 명령의 출처를 암호학적으로 검증하는 것, 즉 MAVLink2 서명 (Signing)이다. 이상탐지·불변식 감시는 서명 위에 얹히는 tripwire(조기 경보)이지 1차 차단선이 아니다.
%     \item \textbf{자율성 그 자체가 방어 계층이자 탐지 불변식이 된다.} ISR 체공기가 위치유지 (LOITER) 중이거나 UGV가 사전계획 미션을 AUTO로 수행 중일 때는 조종사의 스틱 입력이 0이어야 정상이다. 이 상태가 "입력이 0이 아니면 곧 오버라이드 주입"이라는 탐지 가능한 불변식을 만든다.
% \end{enumerate}
방산 자율체계 제어 보안의 핵심 난제는 출처 모호성이다. 인증이 적용되지 않는 MAVLink 운용 구성에서 오토파일럿은 정당한 명령과 경로상에서 주입되는 위조 명령을 메시지 출처만으로 구별하기 어렵다. 위조된 단일 채널 RC 오버라이드 프레임은 작동기 (actuator) 관점에서 정상 조종 입력과 동일한 제어 효과를 낼 수 있다. 따라서 사후 탐지만으로는 이 공격군을 원천적으로 차단하기 어렵고, 다음 두 가지 방어 설계 원리가 필요하다.

\begin{enumerate}[leftmargin=*]
    \item 1차 방어는 명령 출처 인증이어야 한다. 출처 모호성을 프로토콜 계층에서 줄이는 가장 직접적인 방법은 명령의 출처와 무결성을 암호학적으로 검증하는 것이다. 본 연구는 이를 MAVLink v2 서명 기반 메시지 검증으로 다룬다. 이상탐지와 불변식 감시 (invariant monitoring)는 1차 차단선을 대체하는 기능이 아니라, 인증 위에 배치되는 조기 경보 및 방어 보조 계층으로 해석한다.

    \item 자율성은 방어 계층이자 탐지 불변식이 된다. ISR 체공기가 위치유지 상태에 있거나 UGV가 사전 계획된 AUTO 임무를 수행하는 동안에는 정상 운용 정책상 수동 조종 입력이 없어야 하는 구간이 존재한다. 이 구간에서는 ``입력이 0이어야 한다''는 운용 불변식이 형성되며, 비정상 RC 오버라이드 또는 조향 입력은 탐지 가능한 위반 신호가 된다.
\end{enumerate}

\subsection{연구 범위와 정직성 원칙}

% 본 연구의 모든 공격·방어·평가는 ArduPilot SITL (ArduCopter/ArduRover) 환경과 userspace MITM MAVLink relay 위에서 구축·측정되었다. 실기체, 실제 RF·GNSS 송신, 하드웨어 센서, root/kernel netem은 일절 사용하지 않으며, "위성 링크"도 실제 위성이 아니라 고지연·고손실 BLOS 링크 특성을 userspace relay + shaper + flooder로 에뮬레이션한 것이다. 본 연구는 명확히 \textbf{방어적(defensive)} 성격을 가지며, 문서화된 공격을 구축하는 목적은 방어를 시험·강화하는 데 있다. 신규·탐지불가 공격의 발명, 회피(evasion)의 무기화, 운용자 화면 위조는 하지 않는다.

% 본 보고서를 관통하는 원칙은 \textbf{모든 정량 수치를 실제 근거에서만 인용하고 그 경로를 명시하며, 가짜·추정 수치를 만들지 않는다}는 것이다. 근거가 없는 항목은 미측정 또는 설계 수준 (Design-level)으로 정직하게 표기한다. 이 원칙은 세 구분으로 실현된다.

% \begin{itemize}[leftmargin=*]
%     \item \textbf{주입값 (Injected) $\neq$ 창발 결과 (Emergent).} 공격자가 센서·링크에 강제한 입력값 (예: GPS 40 m 변위)은 입력 손잡이일 뿐, 본 연구의 기준 수치는 실제로 기체·임무에 일어난 결과다. 동일한 GPS 40 m 주입이 정지 체공에서는 실제 이탈 0.009 m (EKF innovation 게이팅이 완전 저항)이나, 자율 항로 비행 중에는 실제 crosstrack 18.31 m (회랑 이탈), 지상 rover에서는 34.99 m (거의 1:1)로 전혀 다르게 나타난다 (5장).
%     \item \textbf{탐지 (detection) $\neq$ 무력화 (mitigation).} 탐지 경보가 임계를 넘은 것과, 비행 스택이 자율적으로 모드를 바꿔 페일세이프를 발동한 것은 다르다. 다수 센서 스푸핑 캡처에서 탐지기는 깨끗하게 경보하지만 자율 페일세이프는 발동하지 않고 비행 스택이 오염된 상태로 계속 비행한다.
%     \item \textbf{소표본·잔여 한계의 명시.} 창발 임무 수치 다수는 소표본(n=2--5)이며, AI 평가 계층에는 여전히 LLM이 LLM을 채점하는 잔여가 있고, 일부 방어는 탐지 전용(detect-only)임을 밝힌다.
% \end{itemize}
본 연구의 공격, 방어, 평가는 모두 ArduPilot SITL 환경과 사용자 공간 중간자 (MITM) MAVLink 중계기 (relay) 위에서 수행한다. 실험 플랫폼은 ArduCopter와 ArduRover를 포함한다. 위성/BLOS 링크는 실제 위성을 사용하지 않고, 사용자 공간 relay, shaper, flooder를 통해 고지연, 고손실, 대역폭 경합 등 BLOS C2 데이터링크의 주요 특성을 모사한다.
%
실험 범위는 순수 소프트웨어 시뮬레이션으로 한정한다. 실기체 (Real vehicle), 실제 RF 송수신, GNSS 신호 송신, 하드웨어 센서, 루트 권한 (root privilege), 커널 수준 네트워크 에뮬레이션 (Kernel-level network emulation, netem)은 사용하지 않는다. 따라서 본 연구의 결과는 실전 배치 성능을 직접 주장하지 않고, 방어 로직과 위협 모델의 재현 가능성을 검증하는 시뮬레이션 근거로 해석한다.
%

본 연구는 명확히 \textbf{방어적 (defensive)} 목적을 갖는다. 문헌상 공격군을 구현하는 이유는 공격 능력을 확장하기 위한 것이 아니라, 방어 체계의 탐지, 차단, 복구 성능을 시험하기 위한 것이다. 신규 탐지불가 공격의 발명, 회피 (evasion) 기법의 무기화, 운용자 화면 위조, 실제 RF 또는 GNSS 환경에 대한 개입은 연구 범위에 포함하지 않는다.
%
정량 결과는 실제 근거에서만 인용하고, 해당 근거 경로를 함께 제시한다. 근거가 없는 항목은 추정값으로 채우지 않고, 미측정 (unmeasured) 또는 설계 수준 (design-level)으로 표기한다. 이 원칙은 다음 세 가지 구분으로 구현한다.

\begin{itemize}[leftmargin=*]
    \item 주입값 (Injected value)과 창발 결과 (emergent outcome)를 구분한다. 공격자가 센서 또는 링크에 강제하는 입력값은 공격 강도를 정의하는 손잡이일 뿐이다. 기준 수치는 실제 기체 상태, 임무 궤적, 페일세이프 상태, 회랑 이탈과 같은 창발적 임무 결과다. 예를 들어 동일한 GPS 40 m 주입은 정지 체공 조건에서 0.009 m 수준의 실제 이탈만 유발하지만, 자율 항로 비행 조건에서는 18.31 m의 crosstrack 회랑 이탈로 이어진다. 지상 rover 조건에서는 34.99 m 수준의 이탈이 나타난다.
    \item 탐지 (detection)와 무력화 (mitigation)를 구분한다. 탐지 경보가 임계값을 넘는 것과, 비행 스택이 모드를 전환하거나 페일세이프를 발동하여 공격 효과를 차단하는 것은 다른 결과다. 일부 센서 스푸핑 캡처에서는 탐지기가 명확한 경보를 발생시키더라도 자율 페일세이프가 발동하지 않으며, 비행 스택이 오염된 센서 입력의 영향을 받은 상태로 계속 임무를 수행할 수 있다.
    \item 소표본 조건과 잔여 한계를 명시한다. 창발 임무 수치 다수는 소표본 ($n=2$--$5$)에 기반한다. AI 평가 계층에는 보조 재판정 등 LLM 기반 판단이 일부 남아 있으며, 일부 방어는 실제 예방으로 연결되지 않는 탐지 전용 항목이다.
\end{itemize}


\subsection{핵심 기여}

본 연구의 기여는 네 가지로 요약된다.

% \textbf{기여 1 --- 3개 방산 도메인에 걸친 실측 red/blue 시스템(5장).} UAV·UGV·위성/BLOS 세 도메인에 걸쳐 14개의 문서화된 공격을 실제 운용 맥락 위에서 구현하고, 각각을 구현된 모듈과 SITL 증거에 매핑하였다. 특히 개별 공격의 나열을 넘어 측정된 공격 합성(kill-chain) 2건(KC-A, KC-B)을 실증하고, 실제 사건에 앵커링된 설계 수준 도트린 킬체인 2건(KC-3, KC-4)과 명확히 구분하였다.

% \textbf{기여 2 --- 탐지$\to$차단$\to$복구가 결합된 배치 가능 방어 전략(6장).} 탐지(D2--D8)와 센서 무결성 감시, 미션·파라미터 무결성을 결합하고, 5개의 서로 다른 예방 메커니즘 클래스가 실제로 ASR을 1.0에서 0.0으로 낮춤을 실증하였다. 예방이 배선되지 않은 곳은 정직하게 detect-only로 표기한다.

% \textbf{기여 3 --- ground-truth 도구를 갖춘 실제 LLM 지휘관 아키텍처(7장).} 3개 도메인별 RED 플래너와 단일 공유 BLUE 지휘관으로 구성된 에이전트 아키텍처를 구축하였다. 독립 blind 재판정을 통해 초기의 "도구가 탐지를 가능케 한다"는 주장을 반증하고, 도구 계층의 가치를 일관성·결정성·감사가능성·플레이북 매핑으로 정직하게 재정의하였다.

% \textbf{기여 4 --- LLM에 의존하지 않는 객관적 oracle 평가(8장).} 채점 루프에 LLM이 전혀 개입하지 않는 결정론적 oracle scorer를 도입하여, 주관적 자가채점을 객관 측정으로 대체하였다.

% 마지막으로 본 연구는 한계를 부록으로 미루지 않고 서론에서부터 전면화한다: (i) 모든 결과는 시뮬레이션 전용이다, (ii) 창발 임무 수치는 소표본(n$\leq$5)이다, (iii) 독립 평가는 여전히 LLM이 LLM을 채점하는 잔여를 갖는다, (iv) 일부 방어는 탐지 전용이며 프리벤션이 배선되지 않은 지점에서는 ASR이 불변이다. 이 한계들은 결함의 은폐가 아니라, 어떤 방어가 실제로 작동하고 어떤 것이 설계 수준인지를 정확히 판별할 수 있게 하는 본 연구의 차별화 요소다.


\begin{enumerate}[leftmargin=*, label=\textbf{\arabic*.}]
    \item 세 개 방산 도메인에 걸친 실측 기반 red/blue 시스템.
    UAV, UGV, 위성/BLOS C2 데이터링크의 세 도메인에 걸쳐 14개의 문서화된 공격을 실제 운용 맥락 위에서 구현하고, 각 공격을 구현 모듈과 SITL 증거에 매핑한다. 또한 개별 공격의 나열에 그치지 않고, 실측 공격 합성 (kill chain) 2건을 제시한다. KC-A는 페일세이프 무력화 이후 링크 상실이 치명적 표류로 이어지는 연쇄 효과를 다룬다. KC-B는 GPS와 기압계 복합 기만이 수평 및 수직 오차 성분의 직교합 (quadrature) 관계를 통해 3차원 위치 오차를 증폭하는 현상을 다룬다. 이와 별도로 실제 사건에 앵커링된 설계 수준의 독트린형 킬체인 2건 (KC-3, KC-4)을 제시하며, 실측 기반 킬체인과 설계 수준 킬체인을 명확히 구분한다.
    \item 탐지, 차단, 복구를 결합하는 배치 지향 방어 전략.
    탐지 (D2--D8), 센서 무결성 감시, 미션·파라미터 무결성 검사를 결합하는 다층 방어 체계를 제시한다. 예방 계층에서는 RC 오버라이드 대응 BRAKE 전환, GPS 기만 대응 EKF 소스 전환, MAVLink 명령 주입 대응 MAVLink v2 서명, 재전송·미션·파라미터 변조 대응 MAVLink v2 서명 기반 검증, 대역폭 플러딩 대응 rate-limit를 다룬다. 방어 디스패치 (dispatch) on/off 대조를 통해 이들 예방 항목이 공격성공률을 1.0에서 0.0으로 낮추는지 검증한다. 예방이 실제 SITL 액추에이션에 연결되지 않는 항목은 탐지 전용으로 표기한다.
    \item Ground-truth 도구를 갖춘 실제 LLM 지휘관 아키텍처.
    세 개 도메인별 RED 플래너와 단일 공유 BLUE 지휘관으로 구성되는 에이전트 아키텍처를 제시한다. 탐지는 결정론적 도구가 담당하고, 상관 (correlation), 귀속 (attribution), 대응 계획 (response planning)은 LLM이 담당한다. 독립 blind 재판정은 ``도구가 탐지를 가능하게 한다''는 초기 주장을 ``도구 계층은 일관성, 결정성, 감사가능성, 플레이북 매핑 측면에서 우위를 제공한다''는 주장으로 재정식화한다.
    \item LLM에 의존하지 않는 객관적 오라클 (oracle) 평가.
    채점 루프에 LLM을 포함하지 않는 결정론적 오라클 채점기 (deterministic oracle scorer)를 도입한다. 오라클은 로그, 상태 전이, 공격 주입 시각, 탐지 경보, 방어 발동 여부, 임무 결과를 기준으로 ASR과 방어 효과를 산출한다. 이를 통해 주관적 자가채점이나 LLM 간 상호평가가 아니라, 재현 가능한 규칙 기반 평가로 red/blue 결과를 비교한다.
\end{enumerate}